% Standalone report fragment; requires amsmath, amssymb, booktabs, hyperref.
\section{The PA194AX11Q6 construction and its correctness}
\label{sec:ax11q6}

The new table is a composition of three previously justified changes: the
eleven-entry AX fallthrough interface, exchanged premises in the sixth
quantifier head, and formula coefficient six. The concrete combination is
new relative to the audited snapshot. It is a new machine and a checked
composition, rather than a newly invented proof calculus or architecture.
The state count remains 194. No reduction below 194 follows from this
construction.

\subsection{Machine convention and exact construction}

There is one two-way infinite binary tape, initially zero everywhere. The
head starts at coordinate zero. The dense working controls are
$0,\ldots,193$, the initial control is 0, and the separate terminal label
194 has no outgoing transition. Each working control has two transitions;
every transition writes one binary symbol and moves exactly one square.
The authoritative table is \texttt{machine/machine.tm}. Its SHA-256 is
\begin{quote}\ttfamily\small
9933fe380b5242981255a0b75da6b305fc88df7031e753411a51efaf7dbe0dab.
\end{quote}
The native numbering uses initial control 74 and terminal label 254;
\texttt{labels.json} gives the exact bijection, including the halt label.

Start with the attributed Q-emitter grammar and the ordinary PA194 runtime.
In its structured register source, make these exact changes:
\begin{enumerate}
\item Replace the unique terminal AX \texttt{breakTo 1} by the existing
      nineteen-instruction Push macro \texttt{node49}. This continues into
      the full MP operand reader, including its additional field.
\item Use \texttt{node163=(node6,node54)} instead of
      \texttt{(node6,node55)}. Successful MP returns the conclusion without
      pushing a dummy zero.
\item Replace each four-increment constructor body by two copies of
      \texttt{(inc 1,node60)}, with \texttt{node60=(inc 1,inc 1)}.
      Thus all 61 constructor bodies multiply Cantor pairing by six.
\end{enumerate}
The Q grammar already exchanges the two curried premises of the sixth
head. The builder checks the complete source splice, all unchanged
successors under relocation, and a free-constructor evaluation of the
entire AX emitter. No informal textual replacement of a numeric constant
is used as the source-binding argument.

The raw table has 135 fixed controls and 67 AX-reader controls. Eight
simultaneous contextual folds leave
\[
                         135+67-8=194
\]
working controls. All 32 combinations of old folded control, predecessor
history bit, and current read bit are checked on the final table. Their
local tape identities establish the quotient uniformly in the surrounding
tape. The primitive arithmetic-box rows are unchanged. The added native
reader is
\[
\begin{array}{c|cc}
 &0&1\\\hline
190&0R330&1R330\\
330&0R61&1R99.
\end{array}
\]
The extra fold identifies native control 41 with 99, using the two actions
$0L30$ and $1L2040$. The remaining seven folds are supplied explicitly in
\texttt{folds.json}; the checker validates their incoming-context
hypotheses as well as their local execution identities.

\subsection{The calibrated numerical calculus}

Put $p(a,b)=(a+b)(a+b+1)/2+b$ and define
\[
 I(A,B)=6p(A,B),\quad M(a,b)=6p(a,b)+1,\quad
 U(a,B)=6p(B,a)+2,\quad E(a,b)=6p(a,b)+3.
\]
Write $A\to B$ for $I(A,B)$, $N(B)=I(B,1)$,
$A\wedge B=N(I(A,N(B)))$, and $\exists_a B=N(U(a,N(B)))$.
Zero is truth; one is the designated contradiction. The six ordinary
heads are
\begin{align*}
 h_1(D,H)&=(\exists_D U(D,H))\to H,\\
 h_2(D,A,H)&=U(D,U(A,H))\to U(A,U(D,H)),\\
 h_3(D,A,H)&=M(H,D)\to(E(H,A)\to M(A,D)),\\
 h_4(D,A,H)&=E(H,D)\to(M(A,H)\to M(A,D)),\\
 h_5(D,H)&=\exists_D E(H,D),\\
 h_6^Q(D,A,H)&=U(D,A\to H)\to(U(D,A)\to U(D,H)).
\end{align*}
Also retain
\[
 \rho_8(D,A,H,u)=((H\to u)\to D)\to((D\to H)\to(A\to H)),
 \quad u\le D,
 \qquad \Gamma(B)=I(B,U(B,B)).
\]
The guarded head is given by
\begin{align*}
 V(a,b,c)&=M(a,b)\to U(c,M(a,b)),\\
 P(a,b,c)&=c\to M(c,b),\\
 S(a,b,c)&=U\bigl(b,(M(a,b)\to P(a,b,c))\wedge(M(c,b)\to c)\bigr),\\
 Q(a,b,c)&=U(c,S(a,b,c)\to N(P(a,b,c))),\\
 G(a,b,c)&=N(V(a,b,c)\to Q(a,b,c)),\qquad c\ne a,b.
\end{align*}
The rules are modus ponens and theorem generalization
$B\mapsto U(i,B)$. Call this calculus $\mathcal C_6^Q$.

The inherited separate-head calibration, uniform radix calibration, and
premise-exchange argument give primitive-recursive proof translations,
with correctness provable in PA, establishing
\[
 \mathrm{PA}\vdash
 \bigl(\operatorname{Ref}(\mathcal C_6^Q)
                   \leftrightarrow\operatorname{Ref}(\mathrm{PA})\bigr).
 \tag{1}
\]
The dependence on these results is explicit. The complete inherited
arguments are retained with their attribution; no new proof-assistant
formalization of (1) is claimed. Their applicability here can be checked
at the following exact interfaces.

For the radix change, hereditary recoding sends
$4p(a,b)+r$ to $6p(F(a),F(b))+r$, with $F(0)=0$. The same injection acts
on formula codes and variable labels; it therefore preserves equality
guards and the self-indexed diagonal head. Its primitive-recursive left
inverse establishes injectivity. Since this recoding need not preserve
numerical order, restricted $\rho_8$ leaves are replaced by the inherited
fixed derivation of unrestricted $\rho_8$; they are not naively recoded
as restricted leaves. In the other direction the ordinary finite-bit
interpretation applies to the four constructor residues, interpreting the
unused residues four and five harmlessly. PA supplies the needed finite
binary-comprehension proofs by full induction on each substituted
arithmetic formula.

For premise exchange, the implicational basis derives
$(P\to(Q\to R))\to(Q\to(P\to R))$ and its reverse by a fixed
K/S/MP derivation (equivalently, bracket abstraction of
$\lambda f.\lambda q.\lambda p.f\,p\,q$). Replace each occurrence
of the sixth head by this fixed derivation from the other ordering.
All remaining heads and rules coincide. These are syntactic proof
translations, not an appeal to semantic soundness of PA.

\subsection{Exact commands and their eleven-entry stack}

A command store is $(t,w,d,h)$: visible top, Cantor stack, numerical
offset, and cached head. Set $[a;w]=p(a,w)$, with total projections;
in particular $[0;0]=0$ and both projections of zero equal zero.
Let $D=d+x$ and $H=h+y$. Define the ten generated entries, in stack
order, by
\[
\begin{array}{c|l}
0&G(H,D,c),\text{ or }0\text{ if }c=D\text{ or }c=H\\
1&\Gamma(c)\\
2&h_6^Q(D,c,H)\\
3&h_5(D,H)\\
4&h_4(D,c,H)\\
5&h_3(D,c,H)\\
6&h_2(D,c,H)\\
7&h_1(D,H)\\
8&\rho_8(D,c,H,d)\\
9&U(d,h).
\end{array}
\]
Write them $q_0,\ldots,q_9$. The new AX command has four parameters:
\[
 \mathsf A_{c,x,y,z}(t,w,d,h)
       =(0,[q_0,\ldots,q_9,t;w],z,0).                 \tag{2}
\]
After the extra Push, the full MP reader consumes $z$. It compares
$I(0,z)$ with zero: at $z=0$ it succeeds with conclusion zero, and
at $z>0$ it fails with top zero. Both cases leave the offset $z$.
Thus $z$ cannot be dropped from the scheduler encoding or assumed zero
in the soundness proof.

The other commands remain
\[
 \mathsf S(t,w,d,h)=
 \begin{cases}\mathsf{halt},&t+h=1,\\
 (h,\operatorname{tl}w,d,\operatorname{hd}w),&t+h\ne1,
 \end{cases}
\]
and, with $b=d+x$,
\[
 \mathsf{MP}_x(t,w,d,h)=
 \begin{cases}(b,w,b,0),&I(t,b)=h,\\
 (0,w,b,0),&I(t,b)\ne h.
 \end{cases}                                             \tag{3}
\]
In particular the Shift guard is the actual sum test; it does not halt
at $(t,h)=(1,1)$.

\subsection{All-store proof-word compiler}

Put $A_0=\mathsf A_{0,0,0,0}$ and use semicolons for command
concatenation. Define
\[
 \mathsf{Push}=A_0;\mathsf S^{10};\mathsf{MP}_0,
 \quad\mathsf{Pop}=A_0;\mathsf S^{12},
 \quad\mathsf{Discard}=\mathsf{MP}_0,
 \quad\mathsf{Offset}_i=\mathsf{MP}_i.
\]
Exact substitution in (2)--(3) proves
\begin{align*}
 \mathsf{Push}(t,w,d,h)&=(0,[t;w],0,0),\\
 \mathsf{Pop}(t,[a;w],d,h)&=(t,w,0,a),\\
 \mathsf{Discard}(t,w,0,h)&=(0,w,0,0),\\
 \mathsf{Offset}_i(0,w,0,0)&=(0,w,i,0).
\end{align*}
These words do not halt. The safety calculation uses precisely that
$q_0,\ldots,q_8$ have residue zero modulo six, whereas
$q_9=U(d,h)$ has residue two and is at least two. Pop is safe even
when $t=1$, since the relevant last test contains $q_9\ge2$.
The identities include the zero-pair alias; the stack is not presumed
to be a free list.

For $0\le j\le8$, from $(t,w,0,0)$ the word
\[
 \mathsf{Discard};\mathsf A_{A,D,H,0};
              \mathsf S^{j+2};\mathsf{Pop}^{9-j}
                                                        \tag{4}
\]
returns $(q_j,w,0,0)$. A dominated leaf with $u\le D$ uses
\[
 \mathsf{Discard};\mathsf{Offset}_u;
 \mathsf A_{A,D-u,H,0};\mathsf S^{10};\mathsf{Pop}.
\]
Generalization is implemented by
\[
 \mathsf{Push};\mathsf{Offset}_i;\mathsf S;
                  \mathsf A_{i,0,0,0};\mathsf S^{11},
\]
which sends $(a,w,0,0)$ to $(U(i,a),w,0,0)$, including $a=1$.

At an MP proof node with proofs $e_1$ of $I(a,b)$ and $e_0$ of $a$,
compile the word
\[
 \mathsf{Discard};\mathsf{Push};W_{e_1};\mathsf{Push};
 W_{e_0};\mathsf{Push};\mathsf S;\mathsf S;
                         \mathsf{MP}_b;\mathsf{Pop}.
\]
The initial Push reserves a zero below the two premises. On
$(0,[a,I(a,b),0;w],0,0)$ the displayed final suffix either halts
at the second Shift when $a=1$, or returns $(b,w,0,0)$.
It remains valid at $b=1$ and when zero pairs collapse.

Induction on a finite numerical proof therefore gives: from every
$(t,w,0,0)$, $W_e$ either halts or returns
$(\operatorname{concl}(e),w,0,0)$. In particular, if $e$ proves one,
\[
                            A_0;W_e;\mathsf S             \tag{5}
\]
halts from every carried store $(t,w,d,h)$. The initial AX clears the
visible fields and offset while preserving a finite stack; the final
plain Shift tests $(1,0)$. No reset between scheduler trials is assumed.
The independent free-term checker verifies these identities and all
explicit branch conditions, including the special cases $1+1$ and
$[0;0]=0$.

\subsection{Scheduling, soundness, and transfer to the literal tape}

For a finite command word put $e(\epsilon)=0$ and
$e(c::v)=1+b_c(e(v))$, where
\[
 b_{\mathsf S}(u)=p(0,u),\qquad
 b_{\mathsf{MP}_x}(u)=p(1,p(x,u)),\qquad
 b_{\mathsf A_{c,x,y,z}}(u)=p(c+2,p(y,p(x,p(z,u)))).
\]
The extra AX operand is decoded in the physical source. Each decoded
continuation is a component of $b_c(u)$, and so decreases strictly
relative to the positive continuation $1+b_c(u)$. Arithmetic command
blocks terminate. The source's successive fresh candidates therefore
finish unless one halts. Prefixing (5) by
$\mathsf A_{k,0,0,0}$ gives unbounded candidate codes
$p(k+2,p(0,p(0,p(0,e(\text{word (5)})))))$.
The all-store property of (5) makes this cofinal padding safe.
PA proves that some such candidate is reached unless a halt occurs
earlier. This proves completeness of the source search for refutations.

For soundness, annotate every visible or cached formula and every
canonical Cantor-stack entry by a $\mathcal C_6^Q$ proof. Zero has a
fixed proof; all new AX entries are valid leaves or generalizations of
annotated formulas, and a failed guard emits zero. MP preserves this
invariant using its two premise proofs, and Shift moves the annotations.
If its sum equals one then one visible summand is one, giving a
refutation. The arbitrary numerical offset $z$ carries no formula
annotation. The construction and the soundness induction are primitive
recursive and PA-formalizable.

The literal source has 3,794 primitive sites and 4,901 outcomes, covering
both decrement branches. The construction symbolic engine checks all
outcomes over arbitrary natural primary registers, arbitrary affine
header gap, and an arbitrary nonempty finite spare-register list.
It performs 1,571,333 symbolic operations, with maximum 490 per outcome.
An independent implementation checks all 4,901 outcomes directly against
the literal table and rederives four generic list lemmas from its rows.
It performs 1,149,402 symbolic operations, with maximum 397 and minimum
positive segment-duration constant three. These checks
establish universal finite simulation contracts, not a conclusion from
sampled numerical traces.

Exactly 58,320 literal transitions from blank produce the complete
clean source frame at native control 84, head $-2$, source site 25,
header offset 48, logical registers $(2,0,0,0,0,0,0,0,0)$, and 26 unit spare
registers. The checker compares the whole finite-support tape, with
blank exterior. Every source primitive has a positive finite physical
simulation, so finite concatenation yields source halt implies literal
halt. Conversely, a first literal halt at time $T$ is incompatible with
$T+1$ positive nonhalting source simulations along the deterministic
same run. PA proves this finite argument.

Combining (1), source soundness/completeness, and this exact simulation
gives the ordinary mathematical theorem
\[
 \boxed{\mathrm{PA}\vdash
   \bigl(\neg\operatorname{Halt}_{\mathrm{blank}}(\mathrm{PA194AX11Q6})
                    \leftrightarrow\operatorname{Con}(\mathrm{PA})\bigr).}
\]
This argument relies on the identified inherited PA calibration and
the stated symbolic simulation rules. It is not a newly kernel-checked
Lean/Rocq theorem, an iPRA equivalence, or a proof that PA is consistent.

\subsection{Why the difference is substantive}

Compared with QN6, AX now saves the gate, consumes an extra command
operand, and can set the offset arbitrarily. Even at all-zero parameters
the new stack is $p(0,W)$ rather than the positive old stack $W$, so
the zero-pair alias does not erase the change. Compared with AX11,
all constructor arithmetic and MP matching use coefficient six, and the
sixth emitted head has exchanged premises. For example
$I_4(0,1)=8$ and $I_6(0,1)=12$; neither operation is a control renaming.
The graph census separately checks that the literal table is absent
from the prior corpus under the stated cosmetic symmetries.
